\section{Síntese dos Resultados}
\label{sec:sintese_resultados}

O estudo de impacto dos componentes morfológicos e arquiteturais levou às seguintes definições:

\begin{itemize}
    \item \textbf{A Supremacia Convolucional ($H_1$):} Confirmou-se que as arquiteturas CNN modernas (como DenseNet e ConvNeXt) superam estatisticamente os modelos baseados em \textit{Transformers} (ViT e Swin) no domínio do algodão. Provou-se que o viés indutivo de localidade das convoluções é matematicamente superior à atenção global para o rastreamento de microtexturas e nós de fibra, refutando a tendência de adoção universal de \textit{Transformers} para tarefas visuais industriais estocásticas.
    
    \item \textbf{O Imperativo Fotométrico ($H_2$):} A pesquisa comprovou que o sinal físico de captura atua como o limite superior de desempenho da inteligência artificial. A estabilização da temperatura de cor para a luz branca (6000 K) mitigou o ruído de percepção de amarelamento ($+b$), corrigindo os saltos semânticos interclasses. Conclui-se que a padronização óptica, ditada pela IN 24/2016, é o pré-requisito irrefutável antes de qualquer otimização algorítmica.
    
    \item \textbf{A Obsolescência da Engenharia Manual ($H_3$):} A injeção de características morfológicas explícitas através de bancos de filtros de Gabor falhou em aumentar a capacidade discriminativa das redes. A degradação estatisticamente significante comprovou que, em topologias profundas, o aprendizado de representações \textit{end-to-end} a partir de sinais RGB puros é superior e menos ruidoso do que a \textit{feature engineering} estática.
    
    \item \textbf{O Ponto de Operação Industrial ($H_4$):} Ao transpor a avaliação abstrata para o chão de fábrica, o estudo consolidou a ConvNeXt como o estado da arte prático. Ao entregar o mesmo rigor estatístico preditivo da líder paramétrica (DenseNet), mas operando com uma latência de inferência quase três vezes menor ($6,97$ ms), atendeu-se à demanda corporativa por precisão aliada a baixo custo computacional e viabilidade para \textit{Edge Computing}.
\end{itemize}